%% LyX 2.4.4 created this file.  For more info, see https://www.lyx.org/.
%% Do not edit unless you really know what you are doing.
\documentclass[english]{article}
\usepackage[T1]{fontenc}
\usepackage[latin9]{inputenc}
\usepackage{enumitem}
\usepackage{amsmath}
\usepackage{amsthm}
\usepackage{geometry}
\geometry{verbose,tmargin=1in,bmargin=1in,lmargin=1in,rmargin=1in}
\PassOptionsToPackage{normalem}{ulem}
\usepackage{ulem}

\makeatletter
%%%%%%%%%%%%%%%%%%%%%%%%%%%%%% Textclass specific LaTeX commands.
\numberwithin{equation}{section}
\numberwithin{figure}{section}
\newlength{\lyxlabelwidth}      % auxiliary length 

%%%%%%%%%%%%%%%%%%%%%%%%%%%%%% User specified LaTeX commands.
\usepackage{fancyhdr}
\usepackage{lastpage}
\pagestyle{fancy}
\usepackage{enumitem}
\usepackage{ifthen}
%sets math fonts to be more readable
\everymath=\expandafter{\the\everymath\displaystyle}
%\DeclareMathSizes{10}{10}{10}{10}
\usepackage{multicol}

\usepackage{tikz}
\usetikzlibrary{plotmarks}
\usepackage{pgfplots}
\pgfplotsset{compat=newest}

\usepackage{calc}
\usepackage[nomessages]{fp}

\usepackage{advdate}    % Advancing/saving dates
\usepackage{datetime}   % Dates formatting
\usepackage{datenumber} % Counters for dates
% Added by lyx2lyx
\setlength{\parskip}{\smallskipamount}
\setlength{\parindent}{0pt}

\makeatother

\usepackage{babel}
\begin{document}
\renewcommand{\author}{Dr.\,James Ashe}

\newcommand{\classNum}{137}

\newcommand{\classSec}{B}

\newcommand{\nameType}{Name:}

\newcommand{\examType}{Exam 2}

\renewcommand{\date}{\formatdate{13}{3}{2013}}

%%First page's header%%%%%%%%%%%%%%%%%%%%%%
\fancypagestyle{firststyle} %%header settings for first pagr
{
	\fancyhead[L]{MTH \classNum{}(\classSec{}) \author{}\\
	\textbf{\examType{}} \date{}}
	\fancyhead[C]{\textbf{\nameType}}
	\setlength{\headsep}{14pt}
}
%%%%%%%%%%%%%%%%%%%%%%%%%%%%%%%%%%%%%%%%%%

%%Next page's header%%%%%%%%%%%%%%%%%%%%%%%
\setlist{nolistsep}
\lhead{\classNum{}(\classSec{})} 
\chead{\textbf{\nameType}} 
\cfoot{\thepage{}~of~\pageref{LastPage}} 
\setlength{\headsep}{5pt} 
%%%%%%%%%%%%%%%%%%%%%%%%%%%%%%%%%%%%%%%%%%%

%%Various settings; see the preamble for other settings%%%%%
%\setenumerate{itemsep=.125cm,leftmargin=12pt,itemindent=0pt,labelwidth=0pt}
\setlist{itemsep=.05cm}
\renewcommand{\labelenumi}{\textbf{\arabic{enumi}.}}
\renewcommand{\labelenumii}{\textbf{\alph{enumii}.}}
\thispagestyle{firststyle}
%%%%%%%%%%%%%%%%%%%%%%%%%%%%%%%%%%%%%%%%%%
\newcommand{\xmax}{0}
\newcommand{\xmin}{-\xmax}
\newcommand{\xstep}{0}
\newcommand{\ymax}{0}
\newcommand{\ymin}{-\ymax}
\newcommand{\ystep}{0} 

\textbf{You must show your work to receive credit. }Unless stated
otherwise, all problems are worth 5 points. 

\emph{Determine whether the following relation defines $y$ as a function
of $x$.}
\begin{enumerate}[resume]
\item $x=\left|y+3\right|-1$ \textbf{No. }By the vertical line test.\\
\renewcommand{\xmax}{8}
\renewcommand{\xmin}{-8}
\renewcommand{\xstep}{0} %must be xmin+n
\renewcommand{\ymax}{8}
\renewcommand{\ymin}{-8}
\renewcommand{\ystep}{0} %must be ymin+n
%%%%%%%
\begin{tikzpicture} 
\begin{axis}[height=3cm, 
	width=3cm, 
	axis lines=middle,
	minor x tick num={4},
	minor y tick num={4},
	grid=none,
	xtick={0,50,...,250},
	ytick={0,10,20,...,40},
	%axis line style={-latex}, 
	xmin=\xmin,xmax=\xmax, 
	ymin=\ymin,ymax=\ymax,
	%restrict y to domain=-1:10,
	%no markers, 
	xlabel={$x$},ylabel={$y$}, 
	%every axis x label/.append style={anchor=north}, 
	%every axis y label/.append style={anchor=east}
	]
	
	\addplot[color=blue,mark=*,mark size=1pt] coordinates {(1,-4) (1,0)};
	\addplot[domain=-1:7,thick,samples=100,draw=red,->]{x-1} node[right] {$$};
	\addplot[domain=-1:4,thick,samples=100,draw=red,->]{-x-3} node[right] {$$};
	\draw[blue,dashed] (axis cs:1,\pgfkeysvalueof{/pgfplots/ymin}) -- (axis cs:1,\pgfkeysvalueof{/pgfplots/ymax});
	
\end{axis}
\end{tikzpicture}
\end{enumerate}
\emph{Find the domain and range of the relation. Determine if the
relation is a function.}
\begin{enumerate}[resume]
\item $g=\left\{ \left(2,7\right),\left(30,-4\right),\left(45,7\right),\left(2,5\right)\right\} $

\begin{enumerate}[resume]
\item Domain: $\left\{ 2,30,45\right\} $
\item Range: $\left\{ 7,-4,5\right\} $
\item Is $g$ a function? \textbf{No. }$\left(2,7\right)$ and $\left(2,5\right)$
mean that $g$ pairs 2 with 7 \uline{and} 5. \vfill
\end{enumerate}
\end{enumerate}
\emph{Sketch the graph of the function, then give its domain and range.}\begin{multicols}{2}
\begin{enumerate}[resume]
\item $y=-2x+4$

\begin{enumerate}[resume]
\item Graph:\\
\renewcommand{\xmax}{10} \renewcommand{\xstep}{5}
\renewcommand{\ymax}{\xmax} \renewcommand{\ystep}{5}
%%%%%%%
\begin{tikzpicture} 
	\begin{axis}[
		axis x line=middle, axis y line=middle,     		
		xtick={-\xmax,-\xstep,...,\xmax},
		ytick={-\ymax,-\ystep,...,\ymax},
		minor x tick num={4},
		minor y tick num={4},
		grid=both,
		xlabel={$$},     
		ylabel={$$},
		xmin=-\xmax.25,xmax=\xmax.25,
		ymin=-\ymax.25,ymax=\ymax.25,     		width=5cm,     		height=5cm] 		
		\addplot[domain=\xmin:\xmax,thick,samples=100,draw=red,->]{-2*x+4} node[right] {$$};
		%\addplot[mark=noner,smooth,domain=-1:1]{}; 
	\end{axis}
\end{tikzpicture}
\item Domain: $\left(-\infty,\infty\right)$, i.e., all real numbers
\item Range: $\left(-\infty,\infty\right)$\vfill\columnbreak
\end{enumerate}
\item $f(x)=\begin{cases}
x+2 & \mbox{for }x>1\\
-x-4 & \mbox{for }x\leq1
\end{cases}$

\begin{enumerate}
\item Graph:\\
\renewcommand{\xmax}{10} \renewcommand{\xstep}{5}
\renewcommand{\ymax}{\xmax} \renewcommand{\ystep}{5}
%%%%%%%
\begin{tikzpicture} 
	\begin{axis}[
		axis x line=middle, axis y line=middle,     		
		xtick={-\xmax,-\xstep,...,\xmax},
		ytick={-\ymax,-\ystep,...,\ymax},
		minor x tick num={4},
		minor y tick num={4},
		grid=both,
		xlabel={$$},     
		ylabel={$$},
		xmin=-\xmax.25,xmax=\xmax.25,
		ymin=-\ymax.25,ymax=\ymax.25,     		width=5cm,     		height=5cm] 		
		\addplot[domain=1.2:\xmax,thick,samples=100,draw=red]{x+2} node[right] {$$};
		\addplot[domain=\xmin:1,thick,samples=100,draw=red]{-x-4} node[right] {$$};
		\addplot[color=red,mark=o,solid] coordinates {(1,3)};
		\addplot[color=red,mark=*,solid] coordinates {(1,-5)};
		%\addplot[mark=noner,smooth,domain=-1:1]{}; 
	\end{axis}
\end{tikzpicture}
\item Domain: $\left(-\infty,\infty\right)$
\item Range: $[-5,\infty)$
\end{enumerate}
\end{enumerate}
\end{multicols}

\emph{Use transformation to graph the function. State the domain and
range using interval notation.}\begin{multicols}{2}
\begin{enumerate}[resume]
\item [\textbf{5.}]\setcounter{enumi}{5}$f(x)=-(x+3)^{2}+5$

\begin{enumerate}[resume]
\item Graph:\\
\renewcommand{\xmax}{10} \renewcommand{\xstep}{5}
\renewcommand{\ymax}{\xmax} \renewcommand{\ystep}{5}
%%%%%%%
\begin{tikzpicture} 
	\begin{axis}[
		axis x line=middle, axis y line=middle,     		
		xtick={-\xmax,-\xstep,...,\xmax},
		ytick={-\ymax,-\ystep,...,\ymax},
		minor x tick num={4},
		minor y tick num={4},
		grid=both,
		xlabel={$$},     
		ylabel={$$},
		xmin=-\xmax.25,xmax=\xmax.25,
		ymin=-\ymax.25,ymax=\ymax.25,     		width=5cm,     		height=5cm] 		
		%\addplot[mark=noner,smooth,domain=-1:1]{}; 
		\addplot[domain=\xmin:\xmax,thick,samples=100,draw=red]{-(x+3)^2+5} node[right] {$$};
		\addplot[color=red,mark=*,solid] coordinates {(-3,5)};
	\end{axis}
\end{tikzpicture}
\item Domain: $\left(-\infty,\infty\right)$
\item Range: $(-\infty,5]$\columnbreak
\end{enumerate}
\item $g(x)=-\frac{1}{2}\left|x-1\right|$

\begin{enumerate}
\item Graph:\\
\renewcommand{\xmax}{10} \renewcommand{\xstep}{5}
\renewcommand{\ymax}{\xmax} \renewcommand{\ystep}{5}
%%%%%%%
\begin{tikzpicture} 
	\begin{axis}[
		axis x line=middle, axis y line=middle,     		
		xtick={-\xmax,-\xstep,...,\xmax},
		ytick={-\ymax,-\ystep,...,\ymax},
		minor x tick num={4},
		minor y tick num={4},
		grid=both,
		xlabel={$$},     
		ylabel={$$},
		xmin=-\xmax.25,xmax=\xmax.25,
		ymin=-\ymax.25,ymax=\ymax.25,     		width=5cm,     		height=5cm] 		
		\addplot[mark=noner,smooth,domain=\xmin:\xmax,blue,dashed,thick]{-3};
		\addplot[domain=\xmin:\xmax,thick,samples=100,draw=red]{-1/2*abs(x-1)} node[right] {$$};
		\addplot[color=red,mark=*,solid] coordinates {(1,0)};
		\addplot[color=blue,mark=*,solid] coordinates {(7,-3)};
		\addplot[color=blue,mark=*,solid] coordinates {(-5,-3)};
	\end{axis}
\end{tikzpicture}
\item Is $g(x)$ one-to-one? i.e., does it have an inverse? \textbf{No.
}By the horizontal line test. 
\end{enumerate}
\end{enumerate}
\end{multicols}\vfill\newpage

\emph{Sketch the graph and state the domain and range. Identify any
intervals on which $f(x)$ is increasing or decreasing.}
\begin{enumerate}[resume]
\item [\textbf{7.}]\setcounter{enumi}{7}$f(x)=-\sqrt{x-3}$

\begin{enumerate}[resume]
\item Graph:\\
\renewcommand{\xmax}{10} \renewcommand{\xstep}{5}
\renewcommand{\ymax}{\xmax} \renewcommand{\ystep}{5}
%%%%%%%
\begin{tikzpicture} 
	\begin{axis}[
		axis x line=middle, axis y line=middle,     		
		xtick={-\xmax,-\xstep,...,\xmax},
		ytick={-\ymax,-\ystep,...,\ymax},
		minor x tick num={4},
		minor y tick num={4},
		grid=both,
		xlabel={$$},     
		ylabel={$$},
		xmin=-\xmax.25,xmax=\xmax.25,
		ymin=-\ymax.25,ymax=\ymax.25,     		width=5cm,     		height=5cm] 		
		%\addplot[mark=noner,smooth,domain=-1:1]{};
		\addplot[domain=\xmin:\xmax,thick,samples=100,draw=red]{-sqrt(x-3)} node[right] {$$};
		\addplot[color=red,mark=*,solid] coordinates {(3,0)};
	\end{axis}
\end{tikzpicture}
\item Increasing: $\emptyset$. It does not increase
\item Decreasing: $\left(3,\infty\right)$
\end{enumerate}
\end{enumerate}
\emph{Use the function $h(x)=4x+5$ to complete the following problems.}
\begin{enumerate}[resume]
\item Find $h(4)$. 

$h(4)=4\left(4\right)+5=21$
\item $h(0)=5$. Find $h^{-1}(5)$

$h^{-1}(5)=0$ since $h(0)=5$.
\end{enumerate}
\emph{Use the provide graph to answer the question.}
\begin{enumerate}[resume]
\item The manufacturing cost per vehicle for production of $x$ vehicles
is shown in the graph. Use the graph to find the average rate of change
of the cost per vehicle from 50 to 100 thousand vehicles.

\begin{multicols}{2}\renewcommand{\xmax}{300}
\renewcommand{\xmin}{0}
\renewcommand{\xstep}{0} %must be xmin+n
\renewcommand{\ymax}{40}
\renewcommand{\ymin}{0}
\renewcommand{\ystep}{4} %must be ymin+n
%%%%%%%
\begin{tikzpicture}
	\begin{axis}[
		axis x line=bottom, axis y line=left,     		
		%minor x tick num={0},
		%minor y tick num={0},
		xtick={0,50,...,250},
		ytick={0,10,20,...,40},
		grid=both,,
		%y label style={at={(-0.1,0.65)}},
		xlabel={$x$ Vehicles (thousands)},
		ylabel={Cost each (thousands of \$)},
		xmin=\xmin,xmax=\xmax,
		ymin=\ymin,ymax=\ymax,     		
		width=5.5cm, height=4.5cm] 
		\addplot[mark=noner,smooth,red,thick,domain=0:300]{(10^4*x+5*10^8)/(1000000*x)+5};
		\addplot[color=blue,mark=*,solid] coordinates {(50,15.1)};
		\addplot[color=blue,mark=*,solid] coordinates {(100,10.1)};
		\addplot[mark=noner,smooth,blue,thick,domain=0:300]{20.1-.1*x};
	\end{axis}
\end{tikzpicture}

Using the points $\left(50,15\right)$ and $\left(100,10\right)$,
\\
Average rate of change \\
$\approx\frac{15-10}{50-100}=-\frac{1}{10}$\end{multicols}
\end{enumerate}
\emph{Use the difference quotient, $\frac{f(x+h)-f(x)}{h}$ for $h\neq0$,
to complete and simplify the following problem.}
\begin{enumerate}[resume]
\item $f(x)=7x+3$.

$\frac{\left[7(x+h)+3\right]-\left[7x+3\right]}{h}=\frac{7x+7h-7x-3}{h}=\frac{7h}{h}=7$\vfill\newpage
\end{enumerate}
\emph{Find the inverse of the following function.}
\begin{enumerate}[resume]
\item $f(x)=-\frac{1}{3}x+1$

\begin{eqnarray*}
y & = & -\frac{1}{3}x+1\\
y-1 & = & -\frac{1}{3}x\\
-3(y-1) & = & x\\
f^{-1}(x) & = & -3\left(x-1\right)
\end{eqnarray*}

\end{enumerate}
\emph{Let $f(x)=x+4$ and $g(x)=x^{2}-x$. Find and simplify the expression. }
\begin{enumerate}[resume]
\item $(f+g)(-6)$

$g(-6)=(-6)^{2}-\left(-6\right)=42$ and $f(-6)=(-6)+4=-2$.\\
So, $(f+g)(-6)=-2+42=40$\vfill
\item $\left(g\circ f\right)(2)$

$f(2)=2+4=6$ and $g(6)=36-6=30$.\\
So, $g(f(2))=g(6)=30$\vfill
\end{enumerate}
\emph{For the given functions, find $\left(f\circ g\right)(x)$ and
$\left(g\circ f\right)(x)$, Determine if $g$ is the inverse of $f$. }
\begin{enumerate}[resume]
\item \emph{$f(x)=x^{2}+1$ }and\emph{ $g(x)=\sqrt{x}$.}

\begin{enumerate}[resume]
\item \emph{$\left(f\circ g\right)(x)$}

$f(g(x))=f(\sqrt{x})=(\sqrt{x})^{2}+1=x+1$\vfill
\item \emph{$\left(g\circ f\right)(x)$}

$g(f(x))=g(x^{2}+1)=\sqrt{x^{2}+1}$\vfill
\item Is $g$ the inverse of $f$? \textbf{No. }Either note that $f(g(x))\neq x$
or $g(f(x))\neq x$. 
\end{enumerate}
\end{enumerate}

\end{document}
